% Reusable digit-level conversion elements. Coordinates are in millimetres.
\usetikzlibrary{arrows.meta,calc,fit,decorations.pathreplacing}
\definecolor{CVAccent}{HTML}{26485E}
\tikzset{
 cv digit/.style={inner sep=1pt,minimum height=5mm,font=\large},
 cv bit/.style={inner sep=1pt,minimum width=5mm,minimum height=5mm,font=\normalsize},
 cv arrow/.style={-{Stealth[length=2mm,width=1.3mm]},line width=.65pt,shorten <=1pt,shorten >=1pt},
 cv note/.style={font=\small,align=center,inner sep=1pt},
 cv padding/.style={draw,dashed,line width=.6pt,inner sep=1.6mm},
 cv brace/.style={decorate,decoration={brace,mirror,amplitude=1.6mm},line width=.65pt}
}
\newcommand{\cvDigit}[4]{\node[cv digit,text=CVAccent] (#1) at (#2,#3) {$\mathrm{#4}$};}
% #3 is a comma-separated list of x-coordinate / bit pairs.
\newcommand{\cvBits}[3]{\foreach \x/\v [count=\i] in {#3}{\node[cv bit] (#1-\i) at (\x,#2) {$\v$};}}
\newcommand{\cvMapArrow}[3]{\draw[cv arrow] (#1.south)--(#2,#3);}
% Fixed canvas keeps the radix point in the same page column between stages.
\newcommand{\cvCanvas}[1]{\path[use as bounding box] (-4,5) rectangle (164,#1);}
\newcommand{\cvFourBrace}[1]{%
 \draw[cv brace] (#1-16,-22)--(#1+16,-22);
 \node[cv note] at (#1,-29) {৪ বিট};}
% #1=name, #2=centre x, #3=left bit x, #4=right bit x, #5=octal digit.
\newcommand{\cvOctalGroup}[5]{%
 \draw[cv brace] (#3-4,-22)--(#4+4,-22);
 \cvDigit{#1}{#2}{-40}{#5}
 \draw[cv arrow] (#2,-25)--(#1.north);}
\newcommand{\cvStep}[2]{\par\vspace{2mm}\noindent{\color{CVAccent}\bfseries #1. #2}\par\vspace{2mm}}
% Quotient of a rung is printed as the dividend of the next rung.
\newcommand{\cvDivisionRung}[5]{%
 \node[cv bit] at (#1,#2) {$#3$};
 \node[cv bit] at (#1+10,#2) {$#4$};
 \draw[line width=.65pt] (#1+5,#2+4)--(#1+5,#2-5)--(#1+37,#2-5);
 \node[cv bit] at (#1+32,#2-12) {$\mathrm{#5}$};}
\newcommand{\cvBraceAt}[3]{\draw[cv brace] (#1,#3)--(#2,#3);}
\newcommand{\cvCompactOctal}[5]{%
 \cvBraceAt{#3-3}{#4+3}{-38}
 \cvDigit{#1}{#2}{-51}{#5}
 \draw[cv arrow] (#2,-41)--(#1.north);}
